\chapter{Urinalysis: Protein in Urine}
\section*{What Is This Test?}
Urine protein testing detects and quantifies the amount of protein excreted in the urine. Normally, the glomerular filtration barrier (fenestrated endothelial cells, glomerular basement membrane, and podocyte slit diaphragms) prevents the passage of large proteins (albumin, 66 kDa, and larger molecules) while allowing the filtration of small proteins (<20--30 kDa) that are subsequently reabsorbed in the proximal tubule by the cubilin-megalin receptor complex. Healthy individuals excrete <150 mg of protein per 24 hours (<15 mg/dL on a spot sample), of which albumin constitutes <30 mg/24 hours. Proteinuria is the cardinal sign of glomerular disease and an independent risk factor for the progression of chronic kidney disease and cardiovascular disease \cite{levey2003national,levey2009proteinuria}. The urine dipstick primarily detects albumin using a colorimetric dye-binding method (tetrabromophenol blue); it is insensitive to Bence Jones proteins (immunoglobulin free light chains, which are detected by sulfosalicylic acid [SSA] precipitation test, urine protein electrophoresis, and UPEP). Quantification of proteinuria uses the spot urine protein-to-creatinine ratio (UPCR, mg/g creatinine) or 24-hour urine protein collection. Microalbuminuria (moderately increased albuminuria, 30--300 mg/g creatinine or 30--300 mg/24 hours) is the earliest clinical sign of diabetic nephropathy \cite{toto2004microalbuminuria}. Macroalbuminuria (severely increased albuminuria, >300 mg/g creatinine or >300 mg/24 hours) represents overt nephropathy. Nephrotic-range proteinuria is defined as >3.5 g/24 hours and is a component of nephrotic syndrome (proteinuria + hypoalbuminemia + edema + hyperlipidemia).
\section*{Alternative Names}
Urine Protein; Urine Albumin; Microalbuminuria; Proteinuria; 24-Hour Urine Protein; Spot Urine Protein-Creatinine Ratio; Urine Dipstick Protein
\section*{Principle}
Urine dipstick for protein: tetrabromophenol blue dye-binding at pH 3.0 (buffered); albumin binding produces a color change from yellow to green/blue. Highly specific for albumin; insensitive to Bence Jones proteins and gamma globulins. False-positive: highly alkaline urine (pH >8.0), concentrated urine, gross hematuria, pyuria, contamination with disinfectants (chlorhexidine). False-negative: dilute urine (specific gravity <1.005), non-albumin proteins. SSA precipitation test: 3\% SSA added to urine; turbidity is proportional to protein concentration; detects all proteins (including Bence Jones). Quantitative: spot urine protein-to-creatinine ratio (UPCR, mg/g) correlates with 24-hour urine protein (g/24 hours). 24-hour urine protein collection: creatinine excretion validates completeness (15--20 mg/kg/day in women, 20--25 mg/kg/day in men).
\section*{History}
Proteinuria was recognized as a sign of kidney disease as early as the 1820s, when Richard Bright correlated dropsical patients with albuminous urine and diseased kidneys at autopsy, giving rise to the term ``Bright's disease'' and establishing the clinical value of urine testing. In 1883, simple bedside tests using heat and nitric acid allowed detection of albumin in the urine. The reagent strip (dipstick) method using tetrabromophenol blue was developed in the 1950s, while sulfosalicylic acid precipitation remains a complementary method for detecting non-albumin proteins. The concept of microalbuminuria emerged in the 1980s, when studies showed that small amounts of albumin in the urine of diabetics predicted overt nephropathy, and it became a cornerstone of diabetic kidney disease screening. Spot urine protein-to-creatinine and albumin-to-creatinine ratios were validated in the 1990s and are now preferred over 24-hour collections in most clinical settings.
\section*{Why This Test Is Ordered}
\begin{itemize}
  \item Routine screening during annual examinations, prenatal visits, and diabetic or hypertensive follow-up to detect early kidney disease
  \item Monitoring of patients with chronic kidney disease, diabetes mellitus, or hypertension to stage disease and guide therapy \cite{kasiske1998meta}
  \item Evaluation of a positive screening dipstick to quantify proteinuria and confirm its significance
  \item Assessment of suspected glomerulonephritis, nephrotic syndrome, or tubulointerstitial disease
  \item Detection of microalbuminuria, the earliest sign of diabetic nephropathy, enabling early intervention with ACE inhibitors or ARBs
  \item Evaluation of unexplained edema, foamy urine, or the nephrotic syndrome (proteinuria, hypoalbuminemia, edema, hyperlipidemia)
  \item Detection of Bence Jones proteinuria when multiple myeloma or light chain disease is suspected (negative dipstick with positive SSA test)
\end{itemize}
\section*{Primary vs Secondary Test}
The urine dipstick is a primary screening test for proteinuria. Quantitative measurement (spot UPCR, ACR, or 24-hour urine protein) is the secondary, confirmatory test used to stage proteinuria, guide management, and monitor response to treatment.
\section*{How This Test Is Performed}
For the dipstick test, a reagent strip is dipped into a fresh, well-mixed urine sample (usually a random voided specimen in a clean cup) and read at the designated time; the color change is compared with a chart and graded 0 to 4+. For quantitative testing, a spot urine sample is sent for protein and creatinine, and the protein-to-creatinine ratio (UPCR, in mg/g) is calculated. Alternatively, the patient collects all urine over 24 hours, pooling it in a collection jug kept refrigerated or on ice for measurement of total protein and creatinine; the creatinine content validates the completeness of the collection (15--20 mg/kg/day in women, 20--25 mg/kg/day in men).
\section*{What Preparation Is Needed}
No fasting is required. For a 24-hour collection, the patient is instructed to discard the first morning void, then collect all urine, including the final first morning void of the next day, into the provided container, which should be refrigerated or kept on ice. The patient should avoid strenuous exercise for 24--48 hours before testing, since exercise can cause transient proteinuria, and should report ongoing infections, fever, or medications. Spot specimens should be collected from a fresh void, with first morning urine preferred because it reduces orthostatic and activity-related proteinuria.
\section*{Contraindications and Precautions}
There are no contraindications to urine testing. Precautions include interpreting dipstick results with knowledge of false positives (highly alkaline urine, concentrated urine, hematuria, pyuria, and disinfectant contamination) and false negatives (dilute urine and non-albumin proteins). A dipstick that is negative for protein but positive on the SSA test suggests Bence Jones proteinuria. Orthostatic proteinuria is confirmed by a normal overnight recumbent specimen with protein present in daytime upright specimens.
\section*{Risks}
No significant risks. The test is entirely noninvasive and involves only urine collection.
\section*{What Happens After the Test}
Abnormal results are quantified with a spot UPCR or ACR, or a 24-hour urine protein, and interpreted with serum creatinine, eGFR, and blood pressure. Persistent proteinuria, especially greater than 1 g/24 hours or nephrotic-range proteinuria, prompts further evaluation with serum albumin, lipid profile, urine protein electrophoresis, and possible renal biopsy. Diabetic patients with microalbuminuria are treated with ACE inhibitors or ARBs and followed with repeat ACR measurements.
\section*{Understanding Results}
\subsection*{Negative (Normal)}
Dipstick negative. Spot urine protein <15 mg/dL. 24-hour urine protein <150 mg. UPCR <0.15 g/g (150 mg/g). Microalbuminuria negative (<30 mg/g).
\subsection*{Positive}
Mild proteinuria (150 mg--1 g/24 hours): Orthostatic (postural) proteinuria in children/adolescents (urine collected after overnight recumbency is normal; daytime upright urine shows protein), exercise-induced, fever, UTI, hypertension, early diabetic nephropathy, tubulointerstitial disease (low-molecular-weight proteins). Moderate proteinuria (1--3.5 g/24 hours): Glomerulonephritis (IgA nephropathy, lupus nephritis, post-infectious GN, anti-GBM disease), diabetic nephropathy. Nephrotic-range proteinuria (>3.5 g/24 hours): Nephrotic syndrome. Causes: minimal change disease (children), focal segmental glomerulosclerosis (FSGS), membranous nephropathy, membranoproliferative glomerulonephritis, diabetic nephropathy, and amyloidosis. Bence Jones proteinuria: Multiple myeloma, AL amyloidosis, light chain deposition disease. Negative dipstick protein with positive SSA = Bence Jones protein.
\section*{Normal Results}
Dipstick: negative (<15 mg/dL). 24-hour urine protein: <150 mg. UPCR: <0.15 g/g (150 mg/g). Microalbumin (albumin-to-creatinine ratio, ACR): <30 mg/g (normal), 30--300 mg/g (microalbuminuria), >300 mg/g (macroalbuminuria) \cite{kdigo2012clinical}. Orthostatic proteinuria: urine after overnight recumbency <50 mg/8 hours; upright urine increased.
\section*{Follow-up Tests}
Spot urine protein-to-creatinine ratio (UPCR), 24-hour urine protein, spot urine albumin-to-creatinine ratio (ACR), serum creatinine/eGFR, serum albumin, lipid panel (nephrotic syndrome), urine protein electrophoresis (UPEP) and immunofixation for Bence Jones protein, serum free light chain assay, and renal biopsy for persistent proteinuria >1 g/24 hours or nephrotic syndrome.
\section*{References}
\bibliographystyle{unsrtnat}
\bibliography{references}

\clearpage
